% ============================================================
% §03c — GCP: GERT Capture Paths (Normative)
% Author:  Edith — Spec Editor
% Created: 2026-06-05
% Source of truth: design/gert/grammar/gcp.ebnf (v1.0.0-draft)
% ============================================================
%
% RFC 2119 normative language used throughout:
%   MUST / MUST NOT  — absolute requirements for conforming implementations.
%   SHOULD / SHOULD NOT — recommended but not mandatory.
%   MAY              — optional; permitted but not required.
%
% Grammar production citations take the form: gcp.ebnf §<section>.<subsection>.
% Decision ledger citations take the form: decisions.md "<heading>" <key>.

% Local definition — remove if \TODO is promoted to the main preamble.
\providecommand{\TODO}[1]{\textbf{\textcolor{red}{[TODO:~#1]}}}

% ─────────────────────────────────────────────────────────────────────────────
\chapter{GCP\,---\,GERT Capture Paths}
\label{sec:gcp}
% ─────────────────────────────────────────────────────────────────────────────

\begin{quote}
\noindent\textbf{Informative introduction.}
GCP (GERT Capture Path) is the GERT path language for extracting values
from step outputs and binding them to named variables in the run variable
context.
A GCP expression appears as the value of an entry in a \texttt{capture:}
map in every step type that produces output:
\texttt{cli}, \texttt{tool}, \texttt{http\_call}, \texttt{include},
and event-triggered runbooks.

A GCP expression identifies \emph{what} to extract (the source prefix
and optional traversal path) and implicitly determines \emph{how} to
classify the result (scalar or subtree) for downstream GXL and GIS use.
Captured variables are placed in the run variable context map
(\S\ref{sec:gxl}) and are subsequently available as GDP paths in GXL
boolean expressions (\S\ref{sec:gxl}) and as GIS interpolation targets
(\S\ref{sec:gis}).

The relationship between GCP and GDP (GERT Dotted Path) is as follows:
GDP is the path-traversal sub-grammar shared by GXL expressions and GCP
path suffixes.
A GCP expression is a \emph{source-qualified} path: it names a source
prefix (e.g., \texttt{json}, \texttt{http.body}) and optionally appends
a GDP suffix to navigate into the parsed output.
GDP itself is defined in \S\ref{subsec:gxl:gdp} and is referenced here;
it is not duplicated.

GCP is the normative capture-path grammar for Phase~1 and later.
\end{quote}

% ─────────────────────────────────────────────────────────────────────────────
\section{Normative Reference to the Grammar}
\label{sec:gcp:normref}
% ─────────────────────────────────────────────────────────────────────────────

The file \texttt{design/gert/grammar/gcp.ebnf} (version~1.0.0-draft) is the
authoritative grammar for GCP\@.  This chapter is normative prose \emph{about}
that grammar; it does not supersede it.  Where this prose and the grammar
conflict, \textbf{the grammar wins} and the prose is in error.  Conforming
implementations MUST derive their capture-path parsers from
\texttt{gcp.ebnf}, not from this chapter alone.

GCP values are YAML scalar strings that appear as the right-hand side of
entries in a \texttt{capture:} map.
The GCP parser receives the raw string (after YAML parsing; leading and
trailing whitespace is stripped by the YAML parser before delivery).
The GCP grammar contains no operators, no literals, and no keywords;
it is a pure path language.

(\texttt{gcp.ebnf} \S1, \S2.)

% ─────────────────────────────────────────────────────────────────────────────
\section{Source Prefixes}
\label{sec:gcp:sources}
% ─────────────────────────────────────────────────────────────────────────────

Every GCP expression begins with a \emph{source prefix} that identifies
the origin of the value to be captured.
The top-level GCP production is:

\begin{minted}{text}
CapturePath = StepCapture | HttpCapture | EventCapture | LocalCapture ;
\end{minted}

Conforming parsers MUST reject any \texttt{CapturePath} whose first
segment does not match a recognised source prefix with error
\texttt{GCP-PARSE-001}.

(\texttt{gcp.ebnf} \S3.)

% ─────────────────────────────────────────────────────────────────────────────
\subsection{Local Step Capture}
\label{subsec:gcp:sources:local}

A \emph{local} capture references the current step's own output.
This is the most common form and is used in the \texttt{capture:} blocks
of \texttt{cli}, \texttt{tool}, and \texttt{http\_call} steps.

\begin{minted}{text}
LocalCapture    = LocalScalar | LocalText | LocalStructured ;
LocalScalar     = "exit_code" ;
LocalText       = ( "stdout" | "stderr" ) [ GDP ] ;
LocalStructured = ( "json" | "yaml" ) [ GDP ] ;
\end{minted}

\noindent The local source prefixes are:

\begin{itemize}
  \item \texttt{exit\_code} --- Returns a \texttt{number} (integer in the
    range 0--255 for POSIX exit codes).
    No GDP suffix is permitted; \texttt{GCP-PARSE-002} is raised if a
    \texttt{.} follows \texttt{exit\_code}.

  \item \texttt{stdout} / \texttt{stderr} --- Without a GDP suffix: returns
    the captured text as a \texttt{string} (scalar capture).
    With a GDP suffix: the runtime parses the text as JSON and navigates
    the path, returning a PJVM value (scalar or subtree depending on the
    resolved value).
    This form is equivalent to \texttt{json} with the same GDP suffix.

  \item \texttt{json} / \texttt{yaml} --- Parses the step's
    \texttt{stdout} as the named format and returns a PJVM value.
    Without a GDP suffix (bare-root capture; OI-GCP-02, ratified):
    the entire parsed root document is captured as a PJVM value
    (\texttt{object}, \texttt{array}, or scalar depending on the
    document root type).
    With a GDP suffix: navigates to the specified sub-value.
    If the output is not valid JSON (or YAML), \texttt{GCP-RESOLVE-004}
    is raised.
\end{itemize}

\noindent\textbf{YAML source parity.}
The \texttt{yaml} source applies YAML~1.2 core schema typing:
YAML scalars map to PJVM \texttt{string}, \texttt{number}, \texttt{boolean},
or \texttt{null} per the core schema;
YAML sequences map to \texttt{array};
YAML mappings map to \texttt{object};
YAML anchors and aliases are fully resolved before path traversal.
YAML~1.1-specific types (timestamps, binary) are mapped to \texttt{string}.
Implementations MUST use YAML~1.2 core schema strictly
(\texttt{gcp.ebnf} \S9, OI-GCP-06).

\noindent\textbf{Canonical spelling.}
The GCP grammar standardizes on \texttt{exit\_code} (snake\_case).
The camelCase form \texttt{exitCode} is not a legal GCP token and MUST be
rejected as \texttt{GCP-PARSE-001} (\texttt{gcp.ebnf} \S3.4, OI-GCP-01).

(\texttt{gcp.ebnf} \S3.4.)

% ─────────────────────────────────────────────────────────────────────────────
\subsection{Cross-Step Capture}
\label{subsec:gcp:sources:step}

A \emph{cross-step} capture references the output of a named peer step
by its step identifier.

\begin{minted}{text}
StepCapture = "step" "." StepId "." StepOutput ;
StepId      = IDENT ;
StepOutput  = StepScalar | StepText | StepStructured ;
StepScalar  = "exit_code" ;
StepText    = ( "stdout" | "stderr" ) [ GDP ] ;
StepStructured = ( "json" | "yaml" ) [ GDP ] ;
\end{minted}

The \texttt{StepId} MUST match a step \texttt{id:} value that is
statically present in the runbook.
A \texttt{step.\{id\}.*} path where \texttt{\{id\}} is not found in the
runbook MUST be rejected at plan time with \texttt{GCP-PARSE-006}.
The planner validates against the \emph{statically reachable} step set;
the definition of ``statically reachable'' is specified in
\S\ref{sec:parse-gate} (OPQ-GATE-02 pending).

The source-field semantics for \texttt{StepScalar}, \texttt{StepText},
and \texttt{StepStructured} are identical to those of their
\texttt{LocalCapture} counterparts (\S\ref{subsec:gcp:sources:local}),
with the additional constraint that the referenced step MUST have
completed execution before the capturing step evaluates its
\texttt{capture:} block; otherwise \texttt{GCP-RESOLVE-001} is raised.

(\texttt{gcp.ebnf} \S3.1.)

% ─────────────────────────────────────────────────────────────────────────────
\subsection{HTTP Response Capture}
\label{subsec:gcp:sources:http}

An \emph{HTTP} capture references the response of an \texttt{http\_call}
step.
The \texttt{http.*} source prefix is available only within the
\texttt{capture:} block of the \texttt{http\_call} step itself or within
steps that immediately follow the \texttt{http\_call} step in the
execution flow; cross-step HTTP references MUST use the
\texttt{step.\{id\}.json.*} form instead.

\begin{minted}{text}
HttpCapture = "http" "." HttpField ;
HttpField   = HttpStatus | HttpBody | HttpHeader ;
HttpStatus  = "status" ;
HttpBody    = "body" [ GDP ] ;
HttpHeader  = "headers" "." HEADER_NAME ;
\end{minted}

\noindent The HTTP source fields are:

\begin{itemize}
  \item \texttt{http.status} --- Returns a \texttt{number}
    (the HTTP status code, e.g., 200, 404).
    No GDP suffix is permitted; \texttt{GCP-PARSE-002} is raised if
    \texttt{.} follows \texttt{status}.

  \item \texttt{http.body} --- Without a GDP suffix: returns the raw
    response body as a \texttt{string} (scalar capture).
    With a GDP suffix: the body MUST be valid JSON;
    the runtime navigates the parsed body and returns the PJVM value
    at the specified path.
    If the body is not valid JSON when a GDP suffix is present,
    \texttt{GCP-RESOLVE-004} is raised.

  \item \texttt{http.headers.\{name\}} --- Returns the value of the named
    response header as a \texttt{string}.
    Header name lookup MUST be case-insensitive (RFC~7230 \S3.2);
    the canonical form in the capture path uses lowercase with hyphens
    preserved (e.g., \texttt{content-type}, \texttt{x-request-id}).
    If the header is absent from the response, the capture returns
    \texttt{null} (not \texttt{GCP-RESOLVE-002});
    HTTP headers are optional by nature.
    No GDP suffix is permitted on header values.
\end{itemize}

(\texttt{gcp.ebnf} \S3.2.)

% ─────────────────────────────────────────────────────────────────────────────
\subsection{Event Capture}
\label{subsec:gcp:sources:event}

An \emph{event} capture references the payload of the triggering event
for event-driven runbooks (those whose execution is initiated by an
\texttt{event.filter} trigger).

\begin{minted}{text}
EventCapture = "event" "." EventField ;
EventField   = EventId | EventBody | EventHeader ;
EventId      = "id" ;
EventBody    = "body" [ GDP ] ;
EventHeader  = "headers" "." HEADER_NAME ;
\end{minted}

\noindent The event source fields are:

\begin{itemize}
  \item \texttt{event.id} --- Returns the event's unique identifier as a
    \texttt{string}.
    No GDP suffix is permitted; \texttt{GCP-PARSE-002} is raised if
    \texttt{.} follows \texttt{id}.

  \item \texttt{event.body} --- Without a GDP suffix: returns the raw
    event body as a \texttt{string}.
    With a GDP suffix: the body MUST be valid JSON;
    navigates to the specified sub-value.

  \item \texttt{event.headers.\{name\}} --- Returns the event header value
    as a \texttt{string}.
    Same case-insensitive lookup rule as \texttt{http.headers.*}.
    Returns \texttt{null} if the header is absent.
\end{itemize}

(\texttt{gcp.ebnf} \S3.3.)

% ─────────────────────────────────────────────────────────────────────────────
\subsection{Source Prefix Reference Table}
\label{subsec:gcp:sources:table}

Table~\ref{tab:gcp-sources} is a quick-reference summary of all legal
GCP source prefixes.
All entries are derived from \texttt{gcp.ebnf} \S8 and confirmed against
Don's expression-evaluation audit (\texttt{design/gert/expression-evaluation-audit.md}).

\begin{table}[ht]
\centering
\begin{tabular}{p{5.0cm}p{2.0cm}p{5.0cm}}
\hline
\textbf{Capture Path Pattern} & \textbf{Returns} & \textbf{Applies To} \\
\hline
\multicolumn{3}{l}{\textit{Local captures (current step output)}} \\
\hline
\texttt{exit\_code}            & \texttt{number}      & \texttt{cli}, \texttt{tool} steps \\
\texttt{stdout}                & \texttt{string}      & \texttt{cli}, \texttt{tool} steps \\
\texttt{stdout.\{gdp\}}        & PJVM value           & \texttt{cli}, \texttt{tool} steps \\
\texttt{stderr}                & \texttt{string}      & \texttt{cli}, \texttt{tool} steps \\
\texttt{stderr.\{gdp\}}        & PJVM value           & \texttt{cli}, \texttt{tool} steps \\
\texttt{json}                  & PJVM root            & \texttt{cli}, \texttt{tool} steps \\
\texttt{json.\{gdp\}}          & PJVM value           & \texttt{cli}, \texttt{tool} steps \\
\texttt{yaml}                  & PJVM root            & \texttt{cli}, \texttt{tool} steps \\
\texttt{yaml.\{gdp\}}          & PJVM value           & \texttt{cli}, \texttt{tool} steps \\
\hline
\multicolumn{3}{l}{\textit{Cross-step captures}} \\
\hline
\texttt{step.\{id\}.exit\_code}        & \texttt{number}  & any step type \\
\texttt{step.\{id\}.stdout}            & \texttt{string}  & any step type \\
\texttt{step.\{id\}.stdout.\{gdp\}}    & PJVM value       & any step type \\
\texttt{step.\{id\}.stderr}            & \texttt{string}  & any step type \\
\texttt{step.\{id\}.json}              & PJVM root        & any step type \\
\texttt{step.\{id\}.json.\{gdp\}}      & PJVM value       & any step type \\
\texttt{step.\{id\}.yaml}              & PJVM root        & any step type \\
\texttt{step.\{id\}.yaml.\{gdp\}}      & PJVM value       & any step type \\
\hline
\multicolumn{3}{l}{\textit{HTTP response captures}} \\
\hline
\texttt{http.status}                   & \texttt{number}  & \texttt{http\_call} steps \\
\texttt{http.body}                     & \texttt{string}  & \texttt{http\_call} steps \\
\texttt{http.body.\{gdp\}}             & PJVM value       & \texttt{http\_call} steps \\
\texttt{http.headers.\{name\}}         & \texttt{string}  & \texttt{http\_call} steps \\
\hline
\multicolumn{3}{l}{\textit{Event captures}} \\
\hline
\texttt{event.id}                      & \texttt{string}  & event-triggered \\
\texttt{event.body}                    & \texttt{string}  & event-triggered \\
\texttt{event.body.\{gdp\}}            & PJVM value       & event-triggered \\
\texttt{event.headers.\{name\}}        & \texttt{string}  & event-triggered \\
\hline
\end{tabular}
\caption{GCP source prefix reference.
  Source: \texttt{gcp.ebnf} \S8.
  Bare \texttt{json} and \texttt{yaml} patterns (no GDP suffix) are
  bare-root captures; see \S\ref{sec:gcp:subtree}.}
\label{tab:gcp-sources}
\end{table}

% ─────────────────────────────────────────────────────────────────────────────
\section{Path Traversal}
\label{sec:gcp:traversal}
% ─────────────────────────────────────────────────────────────────────────────

The optional GDP suffix that follows a structured source prefix
(\texttt{json}, \texttt{yaml}, \texttt{stdout} with path, \texttt{body}
with path) is a GERT Dotted Path:

\begin{minted}{text}
GDP         = "." PathSegment { PathSegment } ;
PathSegment = DotSegment | IndexSegment ;
DotSegment  = "." IDENT ;
IndexSegment = "[" INTEGER "]" ;
\end{minted}

\noindent\textbf{Note on notation.}
In GCP context, a GDP suffix begins with an initial \texttt{.} following
the source prefix (e.g., \texttt{json.status.phase}).
In GXL context (\S\ref{subsec:gxl:gdp}), a GDP begins with an
\texttt{IDENT} (the variable name), with no leading \texttt{.}.
The path-traversal algorithm is identical in both contexts; only the
starting point differs.

Path resolution proceeds as follows (source: \texttt{gcp.ebnf} \S4):

\begin{enumerate}
  \item Start at the root of the parsed document (or PJVM value).
  \item For each \texttt{DotSegment} (\texttt{.field}):
    access the named key of the current \texttt{object}.
    If the current value is not an \texttt{object}, raise
    \texttt{GCP-RESOLVE-002}.
    If the key is not present in the object, raise
    \texttt{GCP-RESOLVE-002}.
  \item For each \texttt{IndexSegment} (\texttt{[N]}):
    access the $N$th element (zero-based) of the current \texttt{array}.
    If the current value is not an \texttt{array}, raise
    \texttt{GCP-RESOLVE-003}.
    If $N \geq \text{array length}$, raise
    \texttt{GCP-RESOLVE-003}.
  \item The PJVM value at the final path segment is the captured value.
\end{enumerate}

For the semantics of path traversal within a GXL expression operating
on a previously captured subtree variable, see \S\ref{sec:gxl:semantics:paths}.
GIS dotted paths and GXL GDP paths walk captured subtrees identically.

\noindent Path constraints:

\begin{itemize}
  \item Negative array indices (e.g., \texttt{[-1]}) MUST be rejected at
    parse time with \texttt{GCP-PARSE-003}.

  \item A double-dot or trailing-dot segment (e.g., \texttt{json..field}
    or \texttt{json.field.}) MUST be rejected at parse time with
    \texttt{GCP-PARSE-004}.
\end{itemize}

(\texttt{gcp.ebnf} \S4.)

% ─────────────────────────────────────────────────────────────────────────────
\section{Subtree Capture Semantics}
\label{sec:gcp:subtree}
% ─────────────────────────────────────────────────────────────────────────────

Every captured value is classified at evaluation time into one of two
categories:

\begin{description}
  \item[Scalar capture.]
    The resolved PJVM value is of type \texttt{null}, \texttt{boolean},
    \texttt{number}, or \texttt{string}.
    The scalar value is stored directly in the run variable context map
    under the declared variable name.
    Subsequent GXL expressions access the variable by name;
    GIS templates interpolate it using the coercion rules in
    \S\ref{sec:gis:coercion}.

  \item[Subtree capture.]
    The resolved PJVM value is of type \texttt{array} or \texttt{object}.
    The \textbf{entire} PJVM composite is stored in the run variable context
    map, verbatim, using the PJVM definition in \S\ref{sec:gis:portable-json}.
    The value is \emph{not} serialized to a string at capture time.
    Subsequent GXL GDP paths navigate into the subtree by name or index
    as if it were a top-level document (see \S\ref{sec:gxl:semantics:paths}).
    GIS interpolation of a subtree variable produces the canonical JSON
    serialization of the stored PJVM value (\S\ref{sec:gis:coercion}).
\end{description}

\noindent\textbf{Bare-root capture} (OI-GCP-02, ratified;
decisions.md ``GXL Stream~A open issue ratifications'').
A GCP expression that names a structured source with no GDP suffix
(e.g., \texttt{json}, \texttt{yaml}, \texttt{step.scan.json}) is a
\emph{bare-root capture}.
It captures the entire parsed root document as a PJVM value.
If the root document is an \texttt{object} or \texttt{array},
the result is a subtree capture.
If the root document is a scalar (a single string, number, boolean, or
null value at the top level), the result is a scalar capture.
The capture classification is determined by the actual root type at
evaluation time, not at plan time.

(\texttt{gcp.ebnf} \S5; Q5, resolved.)

% ─────────────────────────────────────────────────────────────────────────────
\subsection{Type Inference at Plan Time}
\label{subsec:gcp:subtree:inference}

The parser/planner SHOULD infer the likely capture type from the path:

\begin{itemize}
  \item \texttt{exit\_code}, \texttt{http.status}, and \texttt{event.id}
    always produce scalar captures.
  \item Paths ending at a field that the conformance corpus declares as
    a known scalar type MAY be inferred as scalar at plan time.
  \item All other paths have unknown type at plan time; the type is
    determined definitively at evaluation time.
\end{itemize}

Type inference is best-effort.
The only normative constraint is that
\texttt{capture.default:} on any path inferred or observed to resolve
to an \texttt{object} or \texttt{array} MUST be rejected
(\S\ref{sec:gcp:default}).

(\texttt{gcp.ebnf} \S5.2.)

% ─────────────────────────────────────────────────────────────────────────────
\subsection{Subtree Use in Iteration}
\label{subsec:gcp:subtree:iteration}

A subtree capture of type \texttt{array} MAY be used as the source of an
\texttt{iterate.over:} field.
The iteration variable receives each element of the array in order;
each element is typed as a PJVM value (scalar or subtree depending on the
element type).

\begin{minted}{yaml}
capture:
  services: json.services          # subtree capture: array
# in a subsequent step:
iterate.over: services
\end{minted}

(\texttt{gcp.ebnf} \S5.1.3.)

% ─────────────────────────────────────────────────────────────────────────────
\section{\texttt{capture.default:} Policy}
\label{sec:gcp:default}
% ─────────────────────────────────────────────────────────────────────────────

\noindent\textbf{Source:} \texttt{gcp.ebnf} \S6; OQ2, ratified.

% ─────────────────────────────────────────────────────────────────────────────
\subsection{Scalar Default Values}
\label{subsec:gcp:default:scalar}

A \texttt{capture.default:} declaration is PERMITTED for scalar captures.
The schema form is:

\begin{minted}{yaml}
capture:
  variable_name: <CapturePath>
capture_defaults:
  variable_name: <scalar_literal>
\end{minted}

Semantics:

\begin{itemize}
  \item If the capture path fails to resolve
    (\texttt{GCP-RESOLVE-002}, \texttt{GCP-RESOLVE-003}),
    or the resolved value is \texttt{null}:
    the declared default value is used as the captured value.
  \item If the capture path resolves to a non-null scalar:
    the resolved value is used and the default is ignored.
\end{itemize}

The default value MUST be of a PJVM scalar type
(\texttt{null}, \texttt{boolean}, \texttt{number}, or \texttt{string}).
A type mismatch between the declared default and the inferred capture type
MUST be rejected at plan time with \texttt{GCP-TYPE-001}.

(\texttt{gcp.ebnf} \S6.2.)

% ─────────────────────────────────────────────────────────────────────────────
\subsection{Subtree Default Prohibition}
\label{subsec:gcp:default:subtree}

A \texttt{capture.default:} declaration MUST NOT be used with a capture
path that resolves to an \texttt{object} or \texttt{array}.
The parser/planner MUST reject such a combination at plan time with error
\texttt{GCP-DEFAULT-SUBTREE}, which is also surfaced as \texttt{PLAN-003}
at the parse gate (\S\ref{sec:parse-gate:error-codes}).

The prohibition exists for three reasons:

\begin{enumerate}
  \item Specifying a full object or array as a YAML default literal is
    verbose and error-prone.
  \item A missing subtree usually signals a structural change in the
    step's output; the consuming step should be conditionally guarded
    rather than silently fed a partial default.
  \item Subtree default values cannot be type-checked at plan time because
    the subtree's field schema is not declared in the runbook.
\end{enumerate}

\noindent\textbf{Correct pattern for optional subtrees.}
Authors MUST guard subtree-consuming steps with a \texttt{when:} clause:

\begin{minted}{yaml}
# REJECTED at plan time — GCP-DEFAULT-SUBTREE / PLAN-003:
capture:
  config: json.optional_config
capture_defaults:
  config: {}

# CORRECT — guard the consuming step instead:
capture:
  config: json.optional_config    # resolves to null if path missing
# In the consuming step:
when: config != null
\end{minted}

When \texttt{capture.default:} is absent and the path resolves to
\texttt{null}, the variable holds \texttt{null} in the context map.
The \texttt{when:} guard in the consuming step prevents evaluation of any
expression that would fail on a \texttt{null} subtree.

If the planner cannot determine the capture type at plan time, the
\texttt{GCP-DEFAULT-SUBTREE} check is deferred to the first evaluation.
If the path resolves to an \texttt{object} or \texttt{array} at runtime
and a default is declared, \texttt{GCP-DEFAULT-SUBTREE} MUST be raised
at that point.

(\texttt{gcp.ebnf} \S6.3.)

% ─────────────────────────────────────────────────────────────────────────────
\section{Error Class Catalog}
\label{sec:gcp:errors}
% ─────────────────────────────────────────────────────────────────────────────

GCP defines three timing categories for errors:
\emph{parse errors} (raised at load/plan time);
\emph{resolution errors} (raised at evaluation time);
and \emph{default-policy and type errors} (raised at plan time or, if type
cannot be inferred, at first evaluation).

A runbook containing a parse error MUST NOT enter execution.
Resolution and type errors cause the affected step to transition to the
\texttt{failed} state.

The master PLAN-* code catalog is owned by \S\ref{sec:parse-gate:error-codes}.
GCP parse errors are surfaced at the parse gate via \texttt{PLAN-001};
the \texttt{GCP-DEFAULT-SUBTREE} error is additionally aliased as
\texttt{PLAN-003}.

Table~\ref{tab:gcp-errors} enumerates every GCP error code.

\begin{table}[ht]
\centering
\begin{tabular}{p{3.8cm}p{2.0cm}p{4.8cm}p{2.4cm}}
\hline
\textbf{Code} & \textbf{Raised By} & \textbf{Description} & \textbf{Remediation} \\
\hline
\multicolumn{4}{l}{\textit{Parse errors (raised at load/plan time)}} \\
\hline
\texttt{GCP-PARSE-001}
  & GCP parser
  & Unknown source prefix: the first path segment does not match any legal source token (\texttt{step}, \texttt{http}, \texttt{event}, \texttt{stdout}, \texttt{stderr}, \texttt{exit\_code}, \texttt{json}, \texttt{yaml}).  Example: \texttt{xml.root.item}.
  & Use a legal source prefix from Table~\ref{tab:gcp-sources}. \\
\texttt{GCP-PARSE-002}
  & GCP parser
  & Invalid path suffix: a GDP suffix appears after a source field that does not permit one (\texttt{exit\_code}, \texttt{http.status}, \texttt{event.id}, any \texttt{headers.*} field).
  & Remove the trailing path segment. \\
\texttt{GCP-PARSE-003}
  & GCP parser
  & Negative array index: \texttt{[-1]} or any negative integer in an \texttt{IndexSegment}.
  & Use a non-negative integer index. \\
\texttt{GCP-PARSE-004}
  & GCP parser
  & Empty path segment: \texttt{..} (double dot) or trailing dot in the capture path.
  & Remove the duplicate or trailing dot. \\
\texttt{GCP-PARSE-005}
  & GCP parser
  & Invalid header name: a character not permitted in \texttt{HEADER\_NAME} (e.g., space, colon, comma) in an \texttt{http.headers.*} or \texttt{event.headers.*} path.
  & Use a valid lowercase HTTP header name (RFC~7230). \\
\texttt{GCP-PARSE-006}
  & Planner
  & Step identifier not found: \texttt{step.\{id\}.*} where \texttt{\{id\}} does not match any step \texttt{id:} in the runbook.  Semantic plan-time error.
  & Correct the step ID or add the missing step. \\
\hline
\multicolumn{4}{l}{\textit{Resolution errors (raised at evaluation time)}} \\
\hline
\texttt{GCP-RESOLVE-001}
  & Evaluator
  & Source not available: the step referenced by \texttt{step.\{id\}.*} has not yet executed, or its output is not available (e.g., the step was skipped by a \texttt{when:} guard).
  & Re-order steps or use a \texttt{when:} guard. \\
\texttt{GCP-RESOLVE-002}
  & Evaluator
  & Path segment not found: a \texttt{DotSegment} references a key absent from the current \texttt{object}, or the current value is not an \texttt{object}.
  & Verify the field name and output schema, or add a default (\S\ref{subsec:gcp:default:scalar}). \\
\texttt{GCP-RESOLVE-003}
  & Evaluator
  & Array index out of bounds: \texttt{IndexSegment} index $\geq$ array length, or the current value is not an \texttt{array}.
  & Guard with a length check or use a scalar path. \\
\texttt{GCP-RESOLVE-004}
  & Evaluator
  & Output not parseable: the step output could not be parsed as the declared format (\texttt{json} or \texttt{yaml}), including empty output when a structured path is given, invalid JSON syntax, or invalid YAML syntax.
  & Ensure the step writes valid JSON or YAML to stdout before a structured capture is attempted. \\
\texttt{GCP-RESOLVE-005}
  & Evaluator
  & Header absent: \texttt{http.headers.*} or \texttt{event.headers.*} references a header not present in the response or event.  Soft resolve: returns \texttt{null}, not a hard error.  An error is raised only if a downstream GXL expression fails on the \texttt{null} value without a guard.
  & Use a \texttt{capture.default:} for the scalar header variable, or guard downstream use with \texttt{when: variable != null}. \\
\hline
\multicolumn{4}{l}{\textit{Default-policy and type errors (plan time, or deferred to first evaluation)}} \\
\hline
\texttt{GCP-DEFAULT-SUBTREE}
  & Planner / Evaluator
  & A \texttt{capture.default:} is declared for a capture path that resolves to an \texttt{object} or \texttt{array} (OQ2, ratified).  Also surfaced as \texttt{PLAN-003} at the parse gate (\S\ref{sec:parse-gate:error-codes}).  Error message MUST include: variable name, capture path, inferred or observed type (if determinable), and the correction hint.
  & Remove the \texttt{capture.default:} and guard the consuming step with \texttt{when: variable != null}. \\
\texttt{GCP-TYPE-001}
  & Planner
  & Default value type mismatch: the declared \texttt{capture.default:} value is not of a scalar type compatible with the capture path's expected type.  Example: \texttt{exit\_code} capture with a string default.
  & Use a default value of the correct scalar type. \\
\hline
\end{tabular}
\caption{GCP error class catalog.
  Source: \texttt{gcp.ebnf} \S7.}
\label{tab:gcp-errors}
\end{table}
